\documentclass[10pt,a4paper]{article}
\usepackage[a4paper,margin=2.2cm]{geometry}
\usepackage[T1]{fontenc}
\usepackage[utf8]{inputenc}
\usepackage{amsmath,amssymb}
\usepackage{booktabs,array,graphicx}
\usepackage{microtype}
\pagestyle{empty}
\setlength{\parindent}{0pt}
\renewcommand{\arraystretch}{1.12}
\newcommand{\celula}[1]{\begin{tabular}[c]{@{}c@{}}#1\end{tabular}}
\newsavebox{\tabcaixa}
\newcommand{\tabajusta}[1]{\sbox{\tabcaixa}{#1}\ifdim\wd\tabcaixa>\linewidth\resizebox{\linewidth}{!}{\usebox{\tabcaixa}}\else\usebox{\tabcaixa}\fi}
\begin{document}

{\large\bfseries Table S2. Constant predictor against calibrated free-space loss by cell}

\medskip
\small
Draws: split draws at $g = 10$~km, $b = 2$~km with at least one valid test node, by cell. Free-space loss (FSPL) is calibrated on the training median of received power plus modelled path loss, over all training nodes (first calibration) or over the valid training nodes only (second calibration). Counts are draws in which FSPL errs less than the constant predictor on the valid test nodes, and draws in which the constant errs less over all test nodes; the last two rows total the sixteen cells, the second one with no buffer ($b = 0$). Medians are mean absolute errors in dB over the valid test nodes.

\medskip
\centering
\tabajusta{\begin{tabular}{lccccccc}
\toprule
Cell & Draws & \celula{Valid \\ \relax nodes: FSPL \\ \relax lower \\ \relax (first \\ \relax calibration)} & \celula{Valid \\ \relax nodes: FSPL \\ \relax lower \\ \relax (second \\ \relax calibration)} & \celula{All nodes: \\ \relax constant \\ \relax lower \\ \relax (first \\ \relax calibration)} & \celula{All nodes: \\ \relax constant \\ \relax lower \\ \relax (second \\ \relax calibration)} & \celula{Constant, \\ \relax valid \\ \relax (median, \\ \relax dB)} & \celula{FSPL (first \\ \relax calibration), \\ \relax valid \\ \relax (median, \\ \relax dB)} \\
\midrule
Bauru Q1 & 59 & 55/59 & 55/59 & 37/59 & 49/59 & 18.101 & 6.210 \\
Bauru Q2 & 56 & 45/56 & 45/56 & 42/56 & 45/56 & 15.080 & 4.906 \\
Bauru Q3 & 56 & 53/56 & 52/56 & 43/56 & 52/56 & 17.826 & 5.615 \\
Bauru Q4 & 56 & 50/56 & 50/56 & 37/56 & 44/56 & 15.492 & 5.406 \\
Campinas Q1 & 59 & 51/59 & 50/59 & 40/59 & 39/59 & 14.114 & 5.067 \\
Campinas Q2 & 59 & 53/59 & 52/59 & 34/59 & 33/59 & 16.940 & 6.258 \\
Campinas Q3 & 57 & 42/57 & 43/57 & 42/57 & 54/57 & 15.767 & 9.719 \\
Campinas Q4 & 58 & 51/58 & 51/58 & 37/58 & 40/58 & 15.512 & 6.661 \\
Lins Q1 & 56 & 56/56 & 56/56 & 44/56 & 55/56 & 18.888 & 6.392 \\
Lins Q2 & 59 & 59/59 & 59/59 & 26/59 & 48/59 & 21.220 & 6.913 \\
Lins Q3 & 48 & 48/48 & 47/48 & 41/48 & 48/48 & 19.481 & 6.912 \\
Lins Q4 & 53 & 53/53 & 52/53 & 39/53 & 50/53 & 19.775 & 5.799 \\
Sorocaba Q1 & 58 & 38/58 & 39/58 & 44/58 & 48/58 & 11.665 & 7.231 \\
Sorocaba Q2 & 59 & 52/59 & 52/59 & 40/59 & 39/59 & 14.749 & 6.283 \\
Sorocaba Q3 & 55 & 43/55 & 43/55 & 43/55 & 44/55 & 20.172 & 8.030 \\
Sorocaba Q4 & 56 & 50/56 & 50/56 & 40/56 & 41/56 & 14.598 & 4.815 \\
\midrule
All cells & 904 & 799/904 (88.4\%) & 796/904 (88.1\%) & 629/904 (69.6\%) & 729/904 (80.6\%) & & \\
\celula{All cells, \\ \relax no buffer} & 938 & 844/938 (90.0\%) & 847/938 (90.3\%) & 691/938 (73.7\%) & 801/938 (85.4\%) & & \\
\bottomrule
\end{tabular}
}

\end{document}
